% !TEX program = xelatex
% Generated by docs/latex/build.py from the sources pinned in sources.json.
\documentclass[11pt,a4paper]{article}
\usepackage[a4paper,margin=25mm]{geometry}
\usepackage{fontspec}
\setmainfont{lmroman10-regular.otf}[BoldFont=lmroman10-bold.otf,ItalicFont=lmroman10-italic.otf,BoldItalicFont=lmroman10-bolditalic.otf]
\setsansfont{lmsans10-regular.otf}[BoldFont=lmsans10-bold.otf,ItalicFont=lmsans10-oblique.otf]
\IfFontExistsTF{DejaVuSansMono.ttf}{\setmonofont{DejaVuSansMono.ttf}[Scale=0.83]}{\setmonofont{lmmono10-regular.otf}[Scale=0.83]}
\usepackage{amsmath,amssymb,mathtools}
\usepackage{graphicx,calc,longtable,booktabs,array}
\usepackage{fvextra}
\usepackage{xurl}
\usepackage[colorlinks=true,linkcolor=blue,urlcolor=blue,citecolor=blue]{hyperref}
\providecommand{\tightlist}{\setlength{\itemsep}{0pt}\setlength{\parskip}{0pt}}
\setlength{\parindent}{0pt}
\setlength{\parskip}{6pt}
\setlength{\emergencystretch}{3em}
\begin{document}
\section{P4 C2 --- Four disjoint congruence classes}\label{p4-c2-four-disjoint-congruence-classes}

\textbf{Solved.} Let \(a_0\pmod{m_0},\ldots,a_3\pmod{m_3}\) be four pairwise disjoint congruence classes of integers, where every modulus \(m_i\) is a positive integer. Then some pair of moduli has greatest common divisor at least \(4\). Repeated moduli and arbitrary integer residues are allowed.

\subsection{Intersection criterion}\label{intersection-criterion}

For positive integers \(m,n\), the classes \(a\pmod m\) and \(b\pmod n\) intersect if and only if \(\gcd(m,n)\) divides \(b-a\). Indeed, any common integer \(z\) gives \(b-a=(z-a)-(z-b)\), a difference of multiples of \(m\) and \(n\), so their gcd divides \(b-a\). Conversely, put \(d=\gcd(m,n)\) and use B\'{e}zout's identity to write \(d=um+vn\). If \(b-a=td\), then

\[
z=a+tum=b-tvn
\]

belongs to both classes. In particular, if two classes are disjoint, their modulus gcd does \textbf{not} divide the difference of their residues. A gcd of \(1\) is therefore impossible. If a pair has gcd \(d\), their residues must be different modulo \(d\).

\subsection{Proof for four classes}\label{proof-for-four-classes}

Suppose, for a contradiction, that all six pairwise modulus gcds are less than \(4\). Each is positive and cannot equal \(1\) by the intersection criterion. Therefore each pair gcd is either \(2\) or \(3\).

If every modulus is even, every pair gcd is divisible by \(2\), so each equals \(2\). The residues of the first three disjoint classes must then be pairwise different modulo \(2\). There are only two residue classes modulo \(2\), which is impossible.

Otherwise choose an odd modulus \(m_i\). Its gcd with each of the other three moduli cannot equal \(2\), so each of those three gcds equals \(3\). Hence \(3\) divides all four moduli, including \(m_i\). Every pair gcd is now divisible by \(3\); because it is less than \(4\), it equals \(3\). Pairwise disjointness therefore makes the four residues pairwise different modulo \(3\). There are only three residue classes modulo \(3\), another contradiction.

The two cases exhaust all possibilities. Thus some pair of moduli has gcd at least \(4\).

The bound is sharp: \(0\pmod4\), \(1\pmod4\), \(2\pmod4\), and \(3\pmod4\) are pairwise disjoint, and every pair of their moduli has gcd exactly \(4\).

\subsection{Formal verification}\label{formal-verification}

\href{https://github.com/mpelteshki/climbing-to-the-frontier/blob/ed55f76eb8c558825bb629fd2cd8dbf46e5f1314/cagent/p4/lean/ProofPursuit/P4/Basic.lean}{Basic.lean} defines actual set disjointness and proves the B\'{e}zout/CRT intersection direction. \href{https://github.com/mpelteshki/climbing-to-the-frontier/blob/ed55f76eb8c558825bb629fd2cd8dbf46e5f1314/cagent/p4/lean/ProofPursuit/P4/C2.lean}{C2.lean} proves \texttt{ProofPursuit.P4.Statement\ 4} by the two cases above. The \href{https://github.com/mpelteshki/climbing-to-the-frontier/blob/ed55f76eb8c558825bb629fd2cd8dbf46e5f1314/cagent/p4/verification.md}{verification record} reports a successful Lean 4.34.1 build and kernel replay of the importing C3 module. No platform submission or organizer acceptance is claimed.
\end{document}
